\documentclass[11pt]{article}
\usepackage[margin=1in]{geometry}
\usepackage{fontspec}
\defaultfontfeatures{Ligatures=TeX}
\setmainfont{Liberation Serif}
\setsansfont{DejaVu Sans}
\setmonofont{DejaVu Sans Mono}
\usepackage{amsmath,amssymb,amsthm,bm}
\usepackage{hyperref}
\usepackage{enumitem}
\usepackage{xcolor}
\setlength{\parskip}{0.55em}
\setlength{\parindent}{0pt}
\hypersetup{colorlinks=true,linkcolor=blue,urlcolor=blue,citecolor=blue}
\title{Theory-mode report: Characteristic polynomials of non-Hermitian random band matrices near the threshold}
\author{Prepared from Scholar Inbox digest dated 2026-04-20}
\date{}

\begin{document}
\maketitle

\section*{Paper information}
\textbf{Title:} Characteristic polynomials of non-Hermitian random band matrices near the threshold

\textbf{Authors:} Mariya Shcherbina and Tatyana Shcherbina

\textbf{arXiv:} \href{https://arxiv.org/abs/2604.15355}{2604.15355}

\textbf{Digest score:} 98 in the Scholar Inbox digest for 2026-04-20.

\section*{Executive summary}
This paper studies a clean local-statistics observable for non-Hermitian random band matrices: the second correlation function of characteristic polynomials. The key scale is the band width $W$ relative to the matrix size $N$. Prior work had already identified two asymptotic regimes: if $W \gg \sqrt{N}$, the normalized correlation behaves like the Ginibre ensemble, while if $1 \ll W \ll \sqrt{N}$, the observable factorizes. What was missing was the true transition window, namely the critical regime $W \asymp \sqrt{N}$.

The main contribution is a precise asymptotic description exactly at that threshold. When $W = \kappa_* \sqrt{N}(1+o(1))$, the normalized correlation no longer collapses to either the Ginibre formula or the factorized one. Instead, it converges to a nontrivial crossover quantity represented by a one-dimensional differential operator $\mathbb{A}_0$ acting on functions on $[-1,1]$. Informally: the paper identifies the critical scaling limit and packages the answer into an operator semigroup formula $(e^{\mathbb{A}_0}\mathbf 1, \mathbf 1)$.

Conceptually, this is valuable because it confirms that non-Hermitian random band matrices exhibit a threshold phenomenon parallel to the more familiar Hermitian story. Technically, the paper extends a supersymmetric transfer-matrix method from the off-critical regimes into the delicate critical window where the leading contributions balance rather than decisively favor one asymptotic phase.

\section*{Problem setup and intuition}
The model is an $N \times N$ non-Hermitian Gaussian random band matrix $H_N = (H_{jk})$ with centered complex Gaussian entries and covariance profile
\[
\mathbb E\, H_{jk} = 0, \qquad \mathbb E\, |H_{jk}|^2 = J_{jk}, \qquad J = (-W^2\Delta + 1)^{-1}.
\]
Here $\Delta$ is the discrete Laplacian on $\{1,\dots,N\}$ with Neumann boundary conditions. Because $J_{jk}$ decays exponentially once $|j-k| \gg W$, the parameter $W$ acts as the effective bandwidth.

The observable of interest is
\[
\Theta(z_1,z_2) = \mathbb E\!\left[\prod_{s=1}^2 \det(H_N-z_s)\det(H_N-z_s)^*\right],
\]
with microscopic spectral shifts
\[
 z_1 = z + \zeta N^{-1/2}, \qquad z_2 = z - \zeta N^{-1/2}, \qquad |z|<1.
\]
This quantity is not literally the two-point eigenvalue correlation function, but it is a robust proxy for local spectral behavior and often captures universality transitions in a cleaner analytic form.

Why should $W \sim \sqrt N$ be critical? Random band matrices mix two competing mechanisms. A narrow band means limited spatial communication across indices, so the system behaves more like weakly coupled local blocks; a wide band mixes strongly enough to resemble a mean-field ensemble. The threshold $W \approx \sqrt N$ is exactly where diffusion across the chain becomes strong enough to compete with the system size. In the non-Hermitian setting, this competition is reflected in the normalized characteristic-polynomial correlation: below threshold it factorizes, above threshold it matches Ginibre, and at threshold it forms a genuine interpolation.

\section*{Main theorem / main result}
The previous paper by the authors established the two-phase picture away from criticality. Roughly, after normalization by diagonal terms,
\[
\frac{\Theta(z_1,z_2)}{\Theta(z_1,z_1)^{1/2}\Theta(z_2,z_2)^{1/2}}
\]
converges to the Ginibre answer when $W \gg \sqrt N$ and to a factorized expression when $W \ll \sqrt N$.

This paper proves the critical-scale theorem. If
\[
W = \kappa_* N^{1/2}(1+o(1)) \qquad (N\to\infty),
\]
then the same normalized correlation converges to
\[
(e^{\mathbb A_0}\mathbf 1,\mathbf 1),
\]
where the differential operator is
\[
\mathbb A_0\varphi(x)
= \frac{1}{8(\kappa_*u_*)^2}\Big((1-x^2)\varphi''(x)-2x\varphi'(x)\Big) + 2x\varphi(x), \qquad x\in[-1,1],
\]
with bounded boundary behavior at $x=\pm 1$, and
\[
 u_* = \sqrt{1-|z|^2}.
\]

A useful way to read this theorem is: the critical limit is neither a scalar closed form of the off-critical type nor an arbitrary model-dependent mess. It is a structured continuum object---a semigroup generated by a second-order operator with a diffusion part and a potential part. The diffusion coefficient depends on the critical proportionality constant $\kappa_*$ and on the bulk spectral position $z$ through $u_*$. So the threshold regime retains universal structure, but with a genuinely new crossover law.

\section*{Proof architecture}
The proof is technical, but the backbone is quite clean. A compact mental model is the following.

\subsection*{1. Supersymmetric integral representation}
The starting point is a supersymmetric (SUSY) representation of the normalized second correlation function. After standard manipulations, the problem is rewritten as the matrix element of a transfer operator:
\[
\widetilde\Theta(z_1,z_2) = (\mathcal K_\zeta^{N-1} g, g).
\]
This converts the asymptotic question into spectral analysis of an integral operator $\mathcal K_\zeta$.

Why this matters: products over matrix sites become powers of one operator, which is exactly the right form for extracting large-$N$ asymptotics.

\subsection*{2. Localization near the dominant manifold}
The next step is to show that the dominant contribution to the integral comes from matrices $Q$ lying close to the manifold where the phase is maximal, essentially
\[
Q^*Q \approx u_*^2 I_2.
\]
This concentration statement lets the authors restrict the operator to a narrow neighborhood of the saddle manifold. Outside that region, the transfer operator norm loses a definite amount, so those contributions die exponentially fast in iteration.

This is the first structural simplification: the full four-complex-dimensional integration problem collapses to a neighborhood of the relevant saddle geometry.

\subsection*{3. Polar / cylinder decomposition}
Inside the dominant region, they decompose matrices as
\[
Q = U\mathcal R,
\]
with $U\in U(2)$ unitary and $\mathcal R>0$ positive. Then $\mathcal R$ is expanded around $u_*I_2$:
\[
\mathcal R = u_*(I_2 + W^{-1/2}R).
\]
This separates radial fluctuations $R$ from angular variables $U$. The transfer kernel factorizes into a Gaussian-type part in $R$ and a structured group kernel in $U$.

Interpretation: radial modes are approximately quadratic and therefore controllable; angular modes carry the subtle long-range information that survives in the critical limit.

\subsection*{4. Quadratic approximation plus controlled corrections}
The authors expand the effective phase around the saddle. The leading term is quadratic in the radial variable, while the dependence on the spectral shift parameter $\zeta$ enters through a perturbation involving the unitary component $U$. Cubic and higher-order terms are shown to be negligible after careful operator estimates.

Away from criticality, these perturbations are either too weak or too strong relative to the main Gaussian contraction, which is why one recovers the simple two-phase laws. At criticality, however, the perturbation and the contraction are of the same effective order, so neither can be discarded.

\subsection*{5. Reduction to an effective one-dimensional operator}
This is the conceptual heart of the paper. After projecting onto the dominant sector of the transfer operator and tracking how the critical scaling survives under repeated composition, the high-dimensional operator problem reduces to an effective scalar differential operator on $[-1,1]$. That operator is precisely $\mathbb A_0$.

The geometry behind the interval $[-1,1]$ comes from the reduced angular dynamics after symmetry reduction. What remains is an effective diffusion with drift/potential, and the limiting correlation is obtained by exponentiating that generator and testing it against the constant function.

\subsection*{6. Matching and error control}
Finally, the paper proves that the original normalized correlation and the effective semigroup expression differ by a vanishing error. This requires several layers of norm estimates: control of the off-saddle region, comparison of the exact transfer operator with its approximations, and stability of the dominant eigenspaces during the reduction steps.

So the theorem is not merely formal asymptotics; it is a rigorous transfer from the original random-matrix observable to the continuum critical generator.

\section*{Why it matters, limitations, and takeaway}
\textbf{Why it matters.} The result fills the last missing regime in this particular observable for non-Hermitian random band matrices. The threshold $W\sim\sqrt N$ is the natural place where localization-type behavior and mean-field behavior compete. Showing that the answer converges to a nontrivial operator crossover is strong evidence that the non-Hermitian band model has a genuine critical theory, not just two disconnected asymptotic phases.

It also strengthens the analogy with the Hermitian random band matrix program, where $\sqrt N$ is the famous one-dimensional critical scale. Even though characteristic polynomial correlations are not the same as full local eigenvalue statistics, they are analytically accessible and often encode the same phase transition structure in a rigorous form.

\textbf{Limitations.} The model is Gaussian and the observable is specialized. So the paper does not directly prove the full eigenvalue point-process transition for general non-Hermitian band matrices. The critical limit is given as an operator formula rather than a simple closed-form expression, which is mathematically sharp but less immediately transparent. Also, because the method is heavily SUSY-transfer-matrix based, extending it to more general entry distributions will require additional universality machinery.

\textbf{Takeaway.} The paper identifies the correct critical object at bandwidth $W\asymp\sqrt N$: the normalized second correlation of characteristic polynomials converges to a semigroup generated by an explicit differential operator. In short, the threshold regime is real, nontrivial, and rigorously tractable. If one cares about the emergence of critical phenomena in non-Hermitian random band models, this is exactly the kind of result that clarifies what the crossover looks like before the full local-statistics theory is available.

\end{document}
